\uuid{OCvf}
\titre{Application linéaire/Polynômes}
\auteur{hanani ; mzari}
\organisation{exo7}
\licence{CC BY-NC-SA 4.0}
\datecreate{2020-01-20}
\niveau{L1}
\module{Algèbre}
\chapitre{Application linéaire}
\sousChapitre{Définition}
\difficulte{2}
\theme{Applications linéaires, Application linéaire, Polynômes}
\exo7qcmid{lille-2.applications-lineaires.004}

\contenu{
\begin{question}
\qclassification{107.01}
\qtags{motcle=Application linéaire/Polynômes}
On note $\Rr_n[X]$ l'espace des polynômes à coefficients réels de degré $\le n$,  $n\in \Nn$. On considère les deux applications suivantes :
$$\begin{array}{rccc}f:&\Rr_3[X]&\to&\Rr\\
& P&\to &P(0)+P'(0)\,  \end{array}   \quad \mbox{et} \quad \begin{array}{rccc}g:&\Rr_3[X]&\to&\Rr_2[X]\\
& P&\to &1+P'+XP'',\end{array}$$
où $P'$ (resp. $P''$) est la dérivée première (resp. seconde) de $P$. Quelles sont les assertions vraies ?
\begin{answers}
\bad{$f(0)=1$}.
\good{$f$ est une application linéaire}.
\good{$g(0)=1$}.
\bad{$g$ est une application linéaire}.
\end{answers}
\begin{explanations} On vérifie que $f$ est linéaire. Par contre, $g$ ne l'est pas, puisque $g(0)=1\neq 0$.
\end{explanations}
\end{question}
}
